% Generated draft: requires mathematical and source review.
\documentclass[11pt]{article}
\usepackage[margin=25mm]{geometry}
\usepackage{fontspec}
\setmainfont{Latin Modern Roman}
\usepackage{amsmath,amssymb,mathtools,unicode-math}
\setmathfont{Latin Modern Math}
\usepackage{enumitem,needspace,xurl,hyperref}
\setlength{\parindent}{0pt}
\setlength{\parskip}{5pt}
\setlength{\emergencystretch}{3em}
\newcommand{\N}{\mathbb{N}}
\newcommand{\Z}{\mathbb{Z}}
\newcommand{\Q}{\mathbb{Q}}
\newcommand{\R}{\mathbb{R}}
\newcommand{\C}{\mathbb{C}}
\newcommand{\F}{\mathbb{F}}
\newcommand{\A}{\mathbb{A}}
\newcommand{\PP}{\mathbb{P}}
\newcommand{\E}{\mathbb{E}}
\newcommand{\eps}{\varepsilon}
\newcommand{\dd}{\,\mathrm{d}}
\begin{document}
\section*{UnsolvedMath ID: 30004606}
\textbf{Source problem code:} catalogue entry
\subsection{Definitions and mathematical statement}
Under the usual convention that a line arrangement is a finite set of distinct lines, let \(\mathcal A=\{L_1,\ldots,L_d\}\) be a complex line arrangement in an ambient plane \(X\). The supplied record does not specify whether \(X=\PP^2(\C)\) or \(X=\C^2\); the definitions below apply after that choice is fixed. Define the intersection set and the multiplicity of an intersection point by
\[
S(\mathcal A)=\bigcup_{1\le i<j\le d}(L_i\cap L_j),\qquad
m_{\mathcal A}(p)=\#\{i:p\in L_i\}\quad(p\in S(\mathcal A)).
\]
For each integer \(r\ge2\), put
\[
t_r(\mathcal A)=\#\{p\in S(\mathcal A):m_{\mathcal A}(p)=r\}.
\]
An ordinary double point is an intersection point of multiplicity exactly two.

\textbf{Classification problem.} Classify complex line arrangements \(\mathcal A\) satisfying
\[
t_2(\mathcal A)=0,
\]
equivalently, satisfying
\[
\forall p\in S(\mathcal A),\qquad m_{\mathcal A}(p)\ge3.
\]
The condition permits intersection points of multiplicity greater than two. No lower bound on \(d\), nonconcurrency hypothesis, or exclusion of arrangements with no intersection points is supplied. The intended ambient plane and the equivalence relation for the classification are unresolved: the record does not say whether the target is a classification of incidence structures or of geometric realizations up to a specified transformation group.
\subsection{Short English statement}
Classify finite arrangements of distinct complex lines for which every intersection point lies on at least three lines. The supplied record leaves the affine or projective setting and the meaning of equivalence in the classification unspecified.
\subsection{Sources}
\begin{itemize}
\item Supplied local UnsolvedMath registry record, permanent ID \(30004606\), title \emph{Complex Line Arrangements without Double Points}, exact source problem code \texttt{catalogue entry}; collection: \emph{Oberwolfach Reports Open Problems}. The record gives proposed year 2021 and names no proposer.
\item Supplied source citation: \emph{Logarithmic Vector Fields and Freeness of Divisors and Arrangements: New perspectives and applications} (2022), Oberwolfach Reports, pp.\ 225--300, DOI \texttt{10.4171/owr/2021/5}. \url{https://doi.org/10.4171/owr/2021/5}. The background describes Oberwolfach Reports as documenting mathematical workshops and their open-problem discussions; the literature list labels this reference \emph{Oberwolfach Report 5/2021}.
\item Supplied literature reference: Hanumanthu--Harbourne, \emph{Journal of Algebraic Combinatorics}. \url{https://doi.org/10.1007/s10801-020-00987-8}. No article title or further bibliographic details are supplied.
\item Supplied literature assessment dated \texttt{2026-08-22}: the classification remains open, and the literature identifies several known types rather than a complete list. This is supplied dated curation derived from the OWR record, not an independently verified assessment. The source expressly requests independent verification of time-sensitive status and cited claims.
\end{itemize}
\end{document}
